\problem{6.26}{Heisenberg momentum of the oscillator}
For the harmonic oscillator
$H=\p^2/(2m)+m\omega^2\x^2/2$,
Example 6.7 finds
\[
x_H(t)=x\cos\omega t+
\frac{p}{m\omega}\sin\omega t.
\]
The Heisenberg equations are
$\dot x_H=p_H/m$ and
$\dot p_H=-m\omega^2x_H$.
Differentiate the displayed
$x_H(t)$ and multiply by $m$:
\begin{equation}
\boxed{p_H(t)=p\cos\omega t
-m\omega x\sin\omega t.}
\end{equation}
At $t=0$ it reduces to $p$; its
time derivative is indeed
\[
\dot p_H=-\omega p\sin\omega t
-m\omega^2x\cos\omega t
=-m\omega^2x_H(t).
\]
These are exactly the \emph{operator}
forms of classical oscillator
trajectories, with $x,p$
replaced by their time-zero
noncommuting operators.

\problem{6.27}{Free-particle Heisenberg position and momentum}
For $H=p^2/(2m)$,
$[H,p]=0$ and
$[x,H]=\ii\hb p/m$.
Since the commutator $[x,H]$
commutes with $H$, the identity
\[
[x,H^n]=
\sum_{j=0}^{n-1}H^j[x,H]H^{n-1-j}
=n\frac{\ii\hb p}{m}H^{n-1}
\]
holds for every positive integer $n$.
Expand $U(t)=\sum_{n\ge0}(-\ii t/\hb)^nH^n/n!$:
\[
[x,U(t)]=
\sum_{n=1}^{\infty}
\frac{(-\ii t/\hb)^n}{n!}
[x,H^n]=\frac{t}{m}p\,U(t).
\]
Using $U^\dagger U=I$,
\begin{align}
x_H(t)&=U^\dagger xU
=x+U^\dagger[x,U]
=x+\frac{t}{m}p,\\
p_H(t)&=U^\dagger pU=p.
\end{align}
Thus
\begin{equation}
\boxed{x_H(t)=x_H(0)+
\frac{t}{m}p_H(0),\qquad
p_H(t)=p_H(0).}
\end{equation}
The mean values follow classical
uniform-velocity motion, while
the operators themselves still obey
$[x_H(t),p_H(t)]=\ii\hb$.

\problem{6.28}{Infinitesimal time-evolution operator}
Taylor-expand the state in
an infinitesimal interval $\delta$:
\[
\Psi(x,t_0+\delta)
=\Psi(x,t_0)+\delta\,\partial_t\Psi(x,t_0)
+O(\delta^2).
\]
The Schrödinger equation gives
$\partial_t\Psi=-(\ii/\hb)H(t_0)\Psi$.
Substitute:
\[
\Psi(x,t_0+\delta)=
\left[I-\frac{\ii\delta}{\hb}H(t_0)\right]
\Psi(x,t_0)+O(\delta^2).
\]
Therefore
\begin{equation}
\boxed{U(t_0+\delta,t_0)
=I-\frac{\ii\delta}{\hb}H(t_0)+O(\delta^2).}
\end{equation}
When $H$ depends on time, this
first-order expression is valid
without requiring $[H(t_1),H(t_2)]=0$.
A finite interval generally
requires time ordering.

\problem{6.29}{Heisenberg equations directly from the evolution operator}
Take a time-independent Hermitian
Hamiltonian and observable $Q$.
With $U(t)=\ee^{-\ii Ht/\hb}$,
$Q_H(t)=U^\dagger(t)QU(t)$.
Differentiate both factors:
\begin{align}
\frac{\dd Q_H}{\dd t}
&=\frac{\ii}{\hb}U^\dagger HQU
-\frac{\ii}{\hb}U^\dagger QHU\\
&=\frac{\ii}{\hb}[H,Q_H].
\end{align}
Equivalently
\begin{equation}
\boxed{\ii\hb\frac{\dd Q_H}{\dd t}
=[Q_H,H].}
\end{equation}
For a one-dimensional
$H=p^2/(2m)+V(x)$, use
$[x,p^2]=2\ii\hb p$ and
$[p,V(x)]=-\ii\hb V'(x)$.
This yields
\begin{equation}
\boxed{\frac{\dd x_H}{\dd t}=
\frac{p_H}{m},\qquad
\frac{\dd p_H}{\dd t}=
-\frac{\partial V(x_H)}{\partial x_H}.}
\end{equation}
The operator derivative of $V$
is interpreted through the
spectral functional calculus.
This is the Heisenberg analogue
of Hamilton's equations.

\problem{6.30}{Propagators for oscillator and free Gaussian packets}
\textbf{(a) Spectral kernel.}
Let $H\psi_n=E_n\psi_n$ and
$\int\psi_n^*(x')\Psi(x',0)\dd x'=c_n$.
The spectral resolution of $U$ is
$\sum_n\ket n\ee^{-\ii E_nt/\hb}\bra n$,
so
\begin{align}
\Psi(x,t)&=\sum_n\psi_n(x)
\ee^{-\ii E_nt/\hb}c_n\\
&=\int K(x,x';t)\Psi(x',0)\dd x',\\
\boxed{K(x,x';t)}&=
\sum_n\psi_n(x)\ee^{-\ii E_nt/\hb}
\psi_n^*(x').
\end{align}
At $t=0$, completeness gives
$K(x,x';0)=\delta(x-x')$.
\textbf{(b) Harmonic oscillator.}
Let $\theta=\omega t$,
$\xi=\sqrt{m\omega/\hb}\,x$,
$\eta=\sqrt{m\omega/\hb}\,x'$.
Insert the normalized Hermite
eigenfunctions
$\psi_n(x)\propto
\ee^{-\xi^2/2}H_n(\xi)/\sqrt{2^nn!}$
and $E_n=\hb\omega(n+\frac12)$
into the spectral sum.
The generating identity given in
the exercise (Mehler's formula),
with $z=\ee^{-\ii\theta}$,
sums it to
\begin{equation}
\boxed{
K_{\rm HO}(x,x';t)=
\sqrt{\frac{m\omega}{2\pi\ii\hb\sin\omega t}}
\exp\!\left\{
\frac{\ii m\omega}{2\hb\sin\omega t}
\big[(x^2+x'^2)\cos\omega t-2xx'\big]
\right\}.}
\end{equation}
This is valid away from focal
times $\sin\omega t=0$, with the
square-root phase chosen by
continuous evolution from $t=0$;
at those times the kernel is a
distribution with a Maslov phase.
\textbf{(c) Evolve an initially displaced Gaussian.}
Let
\[
\Psi(x',0)=\left(\frac{2a}{\pi}\right)^{1/4}
\ee^{-a(x'-x_0)^2}.
\]
Write
$B=m\omega/(2\hb\sin\theta)$,
$A=a-\ii B\cos\theta$, and
$D=2ax_0-2\ii Bx$.
Completing the square in the
integral $\int K_{\rm HO}\Psi_0\dd x'$
uses
\[
\int_{-\infty}^{\infty}
\ee^{-Ax'^2+Dx'}\dd x'
=\sqrt{\frac{\pi}{A}}\,
\ee^{D^2/(4A)},\qquad
\Re A=a>0.
\]
Consequently
\begin{equation}
\boxed{\begin{aligned}
\Psi(x,t)&=
\left(\frac{2a}{\pi}\right)^{1/4}
\sqrt{\frac{m\omega}{2\ii\hb\sin\theta\, A}}\\
&\quad\times\exp\left[
\ii Bx^2\cos\theta-ax_0^2+
\frac{D^2}{4A}\right].
\end{aligned}}
\end{equation}
This compact expression is exact and
includes the phase.
For $a=m\omega/(2\hb)$, its
probability density is a
minimum-width Gaussian centered at
$\boxed{\av{x(t)}=x_0\cos\omega t}$,
the textbook's special coherent-state case.
\textbf{(d) Free-particle kernel.}
Replace the discrete sum by an
integral of continuum momentum
eigenfunctions:
\[
K_0(x,x';t)=\int_{\mathbb R}
\frac{\dd p}{2\pi\hb}
\exp\left[\frac{\ii p(x-x')}{\hb}
-\frac{\ii p^2t}{2m\hb}\right].
\]
Complete the square in $p$
using the Gaussian integral with
the standard $t\to t-\ii0^+$
prescription:
\begin{equation}
\boxed{K_0(x,x';t)=
\sqrt{\frac{m}{2\pi\ii\hb t}}
\exp\left[\frac{\ii m(x-x')^2}{2\hb t}\right].}
\end{equation}
\textbf{(e) Free spreading Gaussian.}
With $\Psi(x,0)=(2a/\pi)^{1/4}\ee^{-ax^2}$,
another Gaussian integral gives
\begin{equation}
\boxed{\Psi(x,t)=
\left(\frac{2a}{\pi}\right)^{1/4}
\frac{1}{\sqrt{1+2\ii\hb at/m}}
\exp\left[-\frac{ax^2}{
1+2\ii\hb at/m}\right].}
\end{equation}
Set $\tau=2\hb at/m$ and take the
absolute square:
\[
|\Psi(x,t)|^2=
\sqrt{\frac{2a}{\pi}}\,
\frac{1}{\sqrt{1+\tau^2}}\,
\exp[-2ax^2/(1+\tau^2)].
\]
Hence
$\boxed{(\Delta x_t)^2=
(1+\tau^2)/(4a)}$,
which verifies normalized free
wavepacket spreading.
\caveat{Although the textbook informally
calls $|K|^2$ a position-transition
probability, a continuum position
eigenstate is not normalizable.
The actual probability density for a
normalized initial state is $|\Psi(x,t)|^2$.}

\problem{6.31}{Translation of a cusp state via spectral resolution}
The normalized cusp function is
$\psi(x)=\sqrt\lambda\,\ee^{-\lambda|x|}$:
\[
\int_{-\infty}^{\infty}|\psi(x)|^2\dd x
=2\lambda\int_0^\infty
\ee^{-2\lambda x}\dd x=1.
\]
The momentum-space transform is
\begin{align}
\phi(p)&=\frac{\sqrt\lambda}{\sqrt{2\pi\hb}}
\int_{\mathbb R}
\ee^{-\lambda|x|}\ee^{-\ii px/\hb}\dd x\\
&=\frac{\sqrt\lambda}{\sqrt{2\pi\hb}}
\left(\frac1{\lambda+\ii p/\hb}
+\frac1{\lambda-\ii p/\hb}\right)\\
&=\boxed{\sqrt{\frac{2\lambda^3}{\pi\hb}}\,
\frac1{\lambda^2+(p/\hb)^2}}.
\end{align}
Now use the spectral representation of
the translation operator:
\[
(\T(a)\psi)(x)=
\frac1{\sqrt{2\pi\hb}}\int_{\mathbb R}
\ee^{\ii p(x-a)/\hb}\phi(p)\dd p.
\]
It is exactly the inverse Fourier
transform of $\phi$ at $x-a$:
\begin{equation}
\boxed{(\T(a)\psi)(x)=
\sqrt\lambda\,\ee^{-\lambda|x-a|}.}
\end{equation}
This proof is valid despite the
undefined ordinary first derivative
at the cusp $x=0$.
\phys{The spectral definition of a
unitary operator is more general
than termwise Taylor expansion
on smooth wavefunctions.}
